% =====================================================================
%  UML-Klassendiagramm der Anwendung BalanceOS (UML 2.5.1)
%  Abgelesen aus dem eingefrorenen Code (app/src/balanceos, Code-Freeze
%  Phase 2). Pakete als UML-Pakete, wichtigste Klassen je Schicht,
%  Funktionsmodule als «module». Typangaben in Python-Schreibweise.
%  Eine Eigenschaft steht entweder als Attribut oder als Assoziation.
%  Export: docs/diagramme/export/klassendiagramm.pdf
% =====================================================================
\documentclass[tikz,border=4pt]{standalone}
\usepackage{iu-diagramm}

\newcommand{\kfont}{\fontsize{7.3pt}{8.6pt}\selectfont\hyphenpenalty=10000\exhyphenpenalty=10000}
\tikzset{
  kl/.style={rectangle split, rectangle split parts=3, draw, fill=white,
             rectangle split part fill={iuFlaeche,white,white},
             rectangle split part align={center,left,left},
             text width=4.1cm, inner sep=2.2pt, font=\kfont, anchor=north west,
             execute at begin node=\centering},
  kl2/.style={kl, rectangle split parts=2, rectangle split part fill={iuFlaeche,white},
              rectangle split part align={center,left}},
  paket/.style={draw, inner sep=0.2cm},
  rolle/.style={font=\fontsize{6.8pt}{7.5pt}\selectfont, inner sep=1.2pt},
  abh/.style={dashed, -{Straight Barb[length=2mm]}},
  ass/.style={-{Straight Barb[length=2mm]}},
  komp/.style={{Diamond[length=2.6mm,width=2mm]}-{Straight Barb[length=2mm]}},
  gen/.style={-{Triangle[open,length=2.6mm,width=2.6mm]}},
}
% UML-Paket: Rahmen um die Knoten #2, Reiter mit dem Paketnamen #1, Name des Rahmens #3
\newcommand{\paket}[3]{%
  \node[paket, fit=#2] (#3) {};
  \node[draw, fill=white, anchor=south west, inner sep=2pt, font=\fontsize{8pt}{9pt}\selectfont\bfseries]
    at (#3.north west) {#1};}

\begin{document}
\begin{tikzpicture}
\def\xa{0}\def\xb{4.9}\def\xc{9.8}\def\xd{14.7}

% =====================================================================
%  praesentation
% =====================================================================
\node[kl] (hf) at (\xa,0) {%
  \textbf{Hauptfenster}\\«QMainWindow»
  \nodepart{second}\raggedright +kontext : AppKontext
  \nodepart{third}\raggedright +zeige\_seite(index : int)\\+neue\_aktivitaet()};
\node[kl] (se) at (\xb,0) {%
  \textbf{Seite}\\«QWidget»
  \nodepart{second}\raggedright +kontext : AppKontext
  \nodepart{third}\raggedright +aktualisieren()};
\node[kl] (zu) at (\xc,0) {%
  \textbf{Zustand}\\«QObject»
  \nodepart{second}\raggedright +daten\_geaendert : Signal\\+woche\_geaendert : Signal\\-\_kw : KW
  \nodepart{third}\raggedright +kw() : KW\\+setze\_kw(kw : KW)};
\node[kl] (ad) at (\xd,0) {%
  \textbf{AktivitaetDialog}\\«QDialog»
  \nodepart{second}\raggedright -\_kontext : AppKontext
  \nodepart{third}\raggedright -\_pruefen() : list[Fehler]\\-\_speichern()};
\node[kl2] (ds) at (\xb,-2.45) {%
  \textbf{DashboardSeite}
  \nodepart{second}\raggedright +aktualisieren()\\+begruendung\_zeigen(hinweis\\\quad : Hinweis)\\+demodaten\_laden()};
\node[notiz, anchor=north west, text width=3.9cm] (sn) at (\xc,-2.75) {%
  Weitere Unterklassen von \textit{Seite}: WocheSeite, AktivitaetenSeite, StatistikSeite,
  ReflexionSeite, KompassSeite, EinstellungenSeite};
\paket{praesentation}{(hf)(se)(zu)(ad)(ds)(sn)}{pPr}

% =====================================================================
%  anwendung
% =====================================================================
\def\ya{-5.7}
\node[kl, text width=5.4cm] (ak) at (\xa,\ya) {%
  \textbf{AppKontext}\\«dataclass»
  \nodepart{second}\raggedright +db : Datenbank\\+heute : Callable[[], date]\\+repos : Repositories\\
  +stammdaten : StammdatenService\\+aktivitaeten : AktivitaetService\\
  +auswertung : AuswertungService\\+export : ExportService\\+daten : DatenService\\
  +reflexion : ReflexionService
  \nodepart{third}\raggedright \underline{+erstellen(pfad : Path | str,}\\\underline{\quad heute) : AppKontext}};
\node[kl] (au) at (\xc,\ya) {%
  \textbf{AuswertungService}
  \nodepart{second}\raggedright -\_regeln : RegelRepository\\-\_db : Datenbank
  \nodepart{third}\raggedright +wochenbewertung(kw : KW)\\\quad : Wochenbewertung\\+hinweise(kw : KW,\\\quad mit\_ausgeblendeten : bool)\\\quad : list[Hinweis]\\+verlauf(bis\_kw : KW, wochen : int)\\\quad : list[Wochenbewertung]};
\node[kl] (as) at (\xd,\ya) {%
  \textbf{AktivitaetService}
  \nodepart{second}\raggedright -\_db : Datenbank\\-\_heute : Callable[[], date]
  \nodepart{third}\raggedright +fuer\_woche(kw : KW)\\\quad : list[Aktivitaet]\\+speichern(a : Aktivitaet)\\\quad : Aktivitaet\\+serie\_anlegen(vorlage : Aktivitaet,\\\quad wochentage, enddatum)\\\quad : Wiederholungsserie\\+loeschen(aktivitaet\_id : int)};
\paket{anwendung}{(ak)(au)(as)}{pAn}

% =====================================================================
%  domaene
% =====================================================================
\def\yd{-11.55}
\node[kl] (lb) at (\xa,\yd) {%
  \textbf{Lebensbereich}
  \nodepart{second}\raggedright +name : str\\+typ : Bereichstyp\\+soll\_stunden : float\\+gewicht : int\\+aktiv : bool
  \nodepart{third}\raggedright +ist\_standard() : bool};
\node[kl2] (mb) at (\xb,\yd) {%
  «module»\\\textbf{bewertung}
  \nodepart{second}\raggedright +bewerte\_woche(kw : KW,\\\quad bereiche, aktivitaeten,\\\quad einstellungen) : Wochenbewertung\\+zielerfuellung(typ, ist, soll)\\\quad : float};
\node[kl2] (mr) at (\xc,\yd) {%
  «module»\\\textbf{regeln}
  \nodepart{second}\raggedright +pruefe\_regeln(k : RegelKontext)\\\quad : list[Hinweis]};
\node[kl2] (mv) at (\xd,\yd) {%
  «module»\\\textbf{validierung}
  \nodepart{second}\raggedright +pruefe\_aktivitaet(a : Aktivitaet,\\\quad bereich, heute) : list[Fehler]\\+pruefe\_lebensbereich(b, alle)\\\quad : list[Fehler]};
\def\ye{-14.6}
\node[kl] (kw) at (\xa,\ye) {%
  \textbf{KW}
  \nodepart{second}\raggedright +jahr : int\\+woche : int
  \nodepart{third}\raggedright +montag() : date\\+verschoben(wochen : int) : KW};
\node[kl] (wb) at (\xb,\ye) {%
  \textbf{Wochenbewertung}
  \nodepart{second}\raggedright +vollstaendigkeit\_tage : int\\+ausreichende\_daten : bool\\+gesamtscore : int | None
  \nodepart{third}\raggedright +bereich(lebensbereich\_id : int)\\\quad : Bereichsbewertung | None};
\node[kl2] (hi) at (\xc,\ye) {%
  \textbf{Hinweis}
  \nodepart{second}\raggedright +regel : str\\+art : Hinweisart\\+text : str\\+begruendung : Begruendung\\+kw : KW};
\node[kl] (av) at (\xd,\ye) {%
  \textbf{Aktivitaet}
  \nodepart{second}\raggedright +titel : str\\+lebensbereich\_id : int\\+datum : date\\+dauer\_minuten : int\\+status : Aktivitaetsstatus
  \nodepart{third}\raggedright +stunden() : float};
\paket{domaene}{(lb)(mb)(mr)(mv)(kw)(wb)(hi)(av)}{pDo}

% =====================================================================
%  persistenz
% =====================================================================
\def\yp{-18.75}
\node[kl] (db) at (\xa,\yp) {%
  \textbf{Datenbank}
  \nodepart{second}\raggedright +pfad : str\\+verbindung : Connection
  \nodepart{third}\raggedright +transaktion()\\+abfrage(sql : str,\\\quad parameter : tuple) : list[Row]\\+ausfuehren(sql : str,\\\quad parameter : tuple) : Cursor};
\node[kl] (rl) at (\xb,\yp) {%
  \textbf{LebensbereichRepository}
  \nodepart{second}\raggedright \mbox{}
  \nodepart{third}\raggedright +alle() : list[Lebensbereich]\\+nach\_id(bereich\_id : int)\\\quad : Lebensbereich | None};
\node[kl] (rr) at (\xc,\yp) {%
  \textbf{RegelRepository}
  \nodepart{second}\raggedright \mbox{}
  \nodepart{third}\raggedright +alle() : list[Regel]\\+aktiv\_setzen(kennung : str,\\\quad aktiv : bool)};
\node[kl] (ra) at (\xd,\yp) {%
  \textbf{AktivitaetRepository}
  \nodepart{second}\raggedright \mbox{}
  \nodepart{third}\raggedright +im\_zeitraum(von : date, bis : date)\\\quad : list[Aktivitaet]\\+speichern(a : Aktivitaet)\\\quad : Aktivitaet};
\paket{persistenz}{(db)(rl)(rr)(ra)}{pPe}

% =====================================================================
%  Beziehungen
% =====================================================================
% praesentation
\path (se.text split) ++(0,0.12) coordinate (h1);
\draw[komp] (hf.east |- h1) -- (se.west |- h1)
  node[rolle, above, pos=0.72] {seiten} node[rolle, below, pos=0.85] {7};
\path (se.text split) ++(0,-0.12) coordinate (h2);
\draw[ass] (se.east |- h2) -- (zu.west |- h2)
  node[rolle, above, pos=0.7] {zustand} node[rolle, below, pos=0.85] {1};
\path (zu.three split) ++(0,-0.28) coordinate (h3);
\draw[ass] (ad.west |- h3) -- (zu.east |- h3)
  node[rolle, above, pos=0.7] {\_zustand} node[rolle, below, pos=0.85] {1};
\draw[gen] (ds.north) -- (se.south);
% praesentation -> anwendung
\draw[ass] ([xshift=1.0cm]ad.south west) -- ([xshift=1.0cm]as.north west)
  node[rolle, right, pos=0.82] {\_service} node[rolle, left, pos=0.82] {1};
% anwendung -> domaene
\path (au.text split) ++(0,-0.35) coordinate (h4);
\draw[abh] (au.west |- h4) -| (mb.north) node[rolle, above, pos=0.25] {«use»};
\draw[abh] ([xshift=-0.6cm]au.south) -- ([xshift=-0.6cm]au.south |- mr.north)
  node[rolle, right, pos=0.5] {«use»};
\draw[abh] ([xshift=1.0cm]as.south west) -- ([xshift=1.0cm]as.south west |- mv.north)
  node[rolle, right, pos=0.5] {«use»};
% domaene intern
\draw[abh] (mb.south) -- (wb.north) node[rolle, right, pos=0.45] {«create»};
\draw[abh] (mr.south) -- (hi.north) node[rolle, right, pos=0.45] {«create»};
\path (wb.text split) ++(0,0.12) coordinate (h5);
\draw[ass] (wb.west |- h5) -- (kw.east |- h5)
  node[rolle, above, pos=0.7] {kw} node[rolle, below, pos=0.85] {1};
% anwendung -> persistenz (Spuren am rechten Rand)
\path (as.text split) ++(0,0.12) coordinate (h6);
\draw[ass] (as.east |- h6) -- ++(0.45,0) coordinate (l1) -- (l1 |- ra.text split)
  node[rolle, sloped, above, pos=0.93] {\_aktivitaeten} -- (ra.east |- ra.text split)
  node[rolle, above, pos=0.5] {1};
\draw[ass] ([xshift=1.2cm]au.south) -- ++(0,-0.25) coordinate (l2) -- (l2 -| pDo.east) -- ++(0.95,0) coordinate (l3)
  -- (l3 |- ra.three split) node[rolle, sloped, above, pos=0.9] {\_aktivitaeten}
  -- (ra.east |- ra.three split) node[rolle, above, pos=0.5] {1};
% Repositories -> Datenbank
\path (db.text split) ++(0,0.12) coordinate (h7);
\draw[ass] (rl.west |- h7) -- (db.east |- h7)
  node[rolle, above, pos=0.7] {\_db} node[rolle, below, pos=0.85] {1};
\path (db.south) ++(0,-0.3) coordinate (t1) ++(0,-0.3) coordinate (t2);
\draw[ass] (rr.south) -- (rr.south |- t1) -| ([xshift=2.8cm]db.south west)
  node[rolle, right, pos=0.93] {1};
\draw[ass] (ra.south) -- (ra.south |- t2) -| ([xshift=1.8cm]db.south west)
  node[rolle, left, pos=0.93] {1};

\end{tikzpicture}
\end{document}
